\section{Zbiór danych i przygotowanie danych}
\label{sec:dane}

Na potrzeby treningu modelu utworzono własny, syntetyczny zbiór danych zawierający polecenia programistyczne oraz odpowiadające im rozwiązania w języku Python. Dane zostały wygenerowane przy użyciu lokalnie uruchamianego dużego modelu językowego za pośrednictwem środowiska Ollama.

Proces tworzenia zbioru przebiegał w trzech etapach. W pierwszej kolejności wygenerowano 500 tematów związanych z programowaniem. Następnie dla każdego tematu przygotowano po 15 różnych poleceń programistycznych. Ostatnim etapem było wygenerowanie wielu rozwiązań dla każdego polecenia.

Każdy rekord zbioru zawiera temat zadania, treść instrukcji oraz wygenerowany kod. Dodatkowo przechowywane są identyfikatory tematu i instrukcji, numer wariantu rozwiązania oraz informacje związane z kontrolą poprawności kodu.

\begin{table}[htbp]
\centering
\caption{Podsumowanie dostępnych plików danych.}
\label{tab:dane-podsumowanie}
\begin{tabular}{lS}
\toprule
\textbf{Artefakt} & {\textbf{Liczba wierszy}} \\
\midrule
\texttt{data/dataset.csv} & 107554 \\
\texttt{data/topics.csv} & 500  \\
\texttt{data/instructions.csv} & 7500  \\
\texttt{data/dataset.csv} & 106262\\
\bottomrule
\end{tabular}
\end{table}



Przed przekazaniem danych do modelu tekst instrukcji oraz kod źródłowy zostały zamienione na sekwencje tokenów. Przy podziale zbioru na część treningową i walidacyjną uwzględniono powiązania między wariantami. Wszystkie rozwiązania wygenerowane dla tej samej instrukcji znajdowały się w tej samej części zbioru. Zapobiega to sytuacji, w której model jest trenowany na jednym wariancie rozwiązania, a następnie oceniany na bardzo podobnym wariancie tego samego zadania.


%
%Plik \texttt{data/dataset.csv} zawiera kolumny \texttt{topic\_id}, \texttt{instruction\_id}, \texttt{topic}, \texttt{instruction}, \texttt{code}, \texttt{valid}, \texttt{variant\_idx}, \texttt{id}, \texttt{code\_file}, \texttt{compile\_valid}, \texttt{validation\_error} oraz \texttt{runtime\_valid}. Do budowy cache aktywnie używane są co najmniej \texttt{instruction} i \texttt{code}, a filtry mogą wykorzystywać \texttt{valid} i \texttt{compile\_valid}. Liczebności pokazane w tabelach~\ref{tab:dane-podsumowanie} i~\ref{tab:dane-jakosc} pochodzą z materiałów dowodowych i bezpośredniego zliczenia plików CSV.
%
%Funkcja \texttt{ensure\_tokenized\_cache()} w \texttt{src/tokenized\_dataset.py} wykonuje filtrowanie, tokenizację oraz zapis macierzy \texttt{.npy}. Filtr \texttt{DATASET\_REQUIRE\_VALID} ogranicza dane do wierszy z \texttt{valid=True}, a \texttt{DATASET\_COMPILE\_FILTER} może wymagać \texttt{compile\_valid=True} albo tylko wykluczać jawne \texttt{False}. W aktualnych metadanych cache występuje tryb \texttt{exclude\_false}.
%
%\input{tables/cache_summary}
%
%Podział na trening i walidację realizuje \texttt{make\_train\_val\_indices()}. Funkcja może używać \texttt{group\_ids}, dzięki czemu warianty tej samej znormalizowanej instrukcji nie powinny trafiać jednocześnie do części treningowej i walidacyjnej. Nie znaleziono zapisanego \texttt{TRAIN\_RESULT\_JSON}, dlatego faktycznych liczebności podziału dla konkretnego przebiegu nie można potwierdzić.
